\documentclass[11pt, letterpaper]{article}

% -----------------------------------------------------------------------------
%  RigbySpace Discrete Dynamics: Rational De Moivre Rotation
%
%  This document presents a discrete, integer–only construction of rotational
%  dynamics.  It adopts the full RigbySpace preamble, including colour
%  definitions, box environments and page formatting, to align with the master
%  reference.  The exposition introduces a rational state pair, defines surge
%  and cross–synchronisation operators and analyses a four–step cycle that
%  emulates Euler's rotation.  We also describe programmable sequences to
%  synthesise various waveforms.  All definitions are given within
%  ontological boxes, and care is taken to avoid continuum concepts.
% -----------------------------------------------------------------------------

\usepackage[top=.8in, bottom=.8in, left=0.5in, right=0.5in]{geometry}
\usepackage{amsmath, amssymb, amsthm, mathtools}
\usepackage{longtable}
\usepackage{booktabs}
\usepackage{array}
\usepackage{tabularx}
\usepackage{xcolor}
\usepackage{fancyhdr}
\usepackage{titlesec}
\usepackage{setspace}
\usepackage[T1]{fontenc}
\usepackage{enumitem}
\usepackage{tocloft}
\usepackage{tcolorbox}
\tcbuselibrary{breakable, skins}

% -----------------------------------------------------------------------------
%  Colour palette matching the RigbySpace master reference
% -----------------------------------------------------------------------------
\definecolor{RSblue}{RGB}{25, 60, 120}
\definecolor{RSgray}{RGB}{90, 90, 90}
\definecolor{RSgreen}{RGB}{20, 110, 60}
\definecolor{RSred}{RGB}{160, 30, 30}
\definecolor{RSamber}{RGB}{180, 110, 0}
\definecolor{RSpurple}{RGB}{90, 30, 120}

% -----------------------------------------------------------------------------
%  Box environments for ontological statements, external verification and
%  structural observations
% -----------------------------------------------------------------------------
\newtcolorbox{rsbox}{
  enhanced, breakable,
  colback=RSblue!4, colframe=RSblue!60,
  title={\textbf{RS Ontological Statement}},
  fonttitle=\small\bfseries\color{white},
  attach boxed title to top left={yshift=-2mm, xshift=4mm},
  boxed title style={colback=RSblue},
  left=4pt, right=4pt, top=6pt, bottom=4pt
}

\newtcolorbox{verifybox}{
  enhanced, breakable,
  colback=RSgreen!4, colframe=RSgreen!50,
  title={\textbf{External Verification (Continuous Lens)}},
  fonttitle=\small\bfseries\color{white},
  attach boxed title to top left={yshift=-2mm, xshift=4mm},
  boxed title style={colback=RSgreen!70!black},
  left=4pt, right=4pt, top=6pt, bottom=4pt
}

\newtcolorbox{observebox}{
  enhanced, breakable,
  colback=RSamber!5, colframe=RSamber!60,
  title={\textbf{Structural Observation}},
  fonttitle=\small\bfseries\color{white},
  attach boxed title to top left={yshift=-2mm, xshift=4mm},
  boxed title style={colback=RSamber!80!black},
  left=4pt, right=4pt, top=6pt, bottom=4pt
}

% -----------------------------------------------------------------------------
%  Page style: header and footer consistent with the master reference
% -----------------------------------------------------------------------------
\pagestyle{fancy}
\fancyhf{}
\fancyhead[L]{\small\textcolor{RSgray}{RigbySpace Discrete Dynamics}}
\fancyhead[R]{\small\textcolor{RSgray}{Rational De Moivre Rotation — Page \thepage}}
\fancyfoot[L]{\small\textcolor{RSgray}{D. Veneziano}}
\fancyfoot[R]{\small\textcolor{RSgray}{Discrete Unification Framework}}
\renewcommand{\headrulewidth}{0.4pt}

% -----------------------------------------------------------------------------
%  Section formatting
% -----------------------------------------------------------------------------
\titleformat{\section}
  {\large\bfseries\color{RSblue}}{\thesection}{1em}{}[\titlerule]
\titleformat{\subsection}
  {\normalsize\bfseries\color{RSblue}}{\thesubsection}{1em}{}
\titleformat{\subsubsection}
  {\normalsize\bfseries\color{RSgray}}{\thesubsubsection}{1em}{}

% Spacing and indentation
\onehalfspacing
\setlength{\parskip}{4pt}
\setlength{\parindent}{0pt}

% Column type for tabular environments
\newcolumntype{L}[1]{>{\raggedright\arraybackslash}p{#1}}

% Common sets and operators
\newcommand{\Z}{\mathbb{Z}}
\newcommand{\lop}{\lambda}
\newcommand{\aop}{\eta}
\newcommand{\psiop}{\psi}

% Theorem styles
\theoremstyle{definition}
\newtheorem{definition}{Definition}[section]
\newtheorem{axiom}{Axiom}[section]
\newtheorem{remark}{Remark}[section]
\newtheorem{postulate}{Postulate}[section]
\newtheorem{observation}{Observation}[section]

\theoremstyle{plain}
\newtheorem{theorem}{Theorem}[section]
\newtheorem{lemma}{Lemma}[section]
\newtheorem{proposition}{Proposition}[section]
\newtheorem{corollary}{Corollary}[section]

% -----------------------------------------------------------------------------
%  Document body
% -----------------------------------------------------------------------------
\begin{document}

% Title page
\begin{titlepage}
  \centering
  \vspace*{2cm}
  {\Huge\bfseries\color{RSblue} RigbySpace Discrete Dynamics\par}
  \vspace{0.5cm}
  {\LARGE Rational De Moivre Rotation\par}
  \vspace{0.5cm}
  {\large Discrete Emulation of Euler's Identity via Integer Pair Dynamics\par}
  \vspace{2cm}
  \rule{0.6\textwidth}{0.4pt}\par
  \vspace{1cm}
  {\large D. Veneziano\par}
  \vspace{0.8cm}
  {\normalsize RigbySpace Study\par}
  \vspace{1.5cm}
  \rule{0.6\textwidth}{0.4pt}\par
  \vspace{0.5cm}
  {\small This study constructs a rational analogue of De Moivre's rotation
  using only integer arithmetic.  Two streams of unreduced rational pairs
  interact through surge and twist operations to emulate rotation and
  oscillation without appealing to complex numbers or the continuum.\par}
  \vfill
  {\small\color{RSgray}
  RigbySpace is a formally consistent candidate theory internal to its own
  axioms.  No empirical equivalence to external physical theory is asserted
  unless an explicit mapping is provided and derived under that mapping.\par}
\end{titlepage}

\tableofcontents
\newpage

% -----------------------------------------------------------------------------
%  Section 1: Introduction
% -----------------------------------------------------------------------------
\section{Introduction}

In continuous mathematics, complex exponentials and trigonometric functions
describe rotations and oscillations elegantly via Euler's identity.  Within
RigbySpace, however, continuum objects and complex numbers are forbidden.  To
recreate rotational and oscillatory dynamics, we must rely on discrete
operations on unreduced rational pairs.  This document introduces a coupled
system of two integer streams that evolve according to surge and
cross–synchronisation operations.  By alternating growth and exchange, the
system traces a discrete analogue of a rotating vector with damping.  The
construction illustrates how familiar physical behaviours can emerge from
integer arithmetic alone.

\subsection{Motivation and Structural History}

Previous work on the triadic rational process demonstrates how irrational
limits arise from unilateral recursions, yet such systems converge to static
values.  To model time–dependent phenomena—oscillations, waves and rotations—we
introduce coupling between two streams.  Each denominator records the full
history of operations applied to its register; this ``structural history''
serves as a proxy for inertia.  When denominators are exchanged between
registers, the accumulated history migrates, producing a quarter–turn in the
effective plane.  We occasionally use the term ``energy'' to convey this
intuitive picture; it should be understood as a metaphor for structural
history rather than a physical quantity.

% -----------------------------------------------------------------------------
%  Section 2: Discrete Rotation Framework
% -----------------------------------------------------------------------------
\section{Discrete Rotation Framework}

We formalise the components of the system: the rational state pair, the surge
operators that grow each register and the cross–synchronisation operator that
exchanges structural history.  All definitions appear in ontological boxes.

\subsection{Rational State Pair}

\begin{rsbox}
  \begin{definition}[Rational State Pair]
  A \emph{rational state} at tick $n$ is an ordered pair of positive integer
  vectors $S_n = (A_n,B_n)$ with
  \begin{align*}
    A_n &= (n_a,d_a) \in \Z_{>0}^2, \\
    B_n &= (n_b,d_b) \in \Z_{>0}^2.
  \end{align*}
  The 
  quantities $v_A = n_a/d_a$ and $v_B = n_b/d_b$ are the rational values
  represented by the unreduced pairs.  No cancellation of factors occurs; the
  denominators $d_a$ and $d_b$ record the entire structural history of
  operations.  The register $A_n$ is designated the ``real'' component and
  $B_n$ the ``imaginary'' component.  Unless otherwise noted, we assume
  initial conditions $n_a,d_a,n_b,d_b \in \Z_{>0}$, with at least one
  denominator strictly greater than one to avoid triviality.
  \end{definition}
\end{rsbox}

\subsection{Surge Operators}

Growth of a register is effected by two surge operators analogous to those in
the master reference.

\begin{rsbox}
  \begin{definition}[Aggregate Surge]
  The \emph{aggregate surge} $\aop$ acts on a component $(n,d)$ by
  \begin{equation*}
    \aop(n,d) = (n + d, n).
  \end{equation*}
  This operation increases the numerator by the denominator and replaces the
  denominator with the previous numerator.  In terms of rational values,
  $v\mapsto 1 + 1/v$.  Repeated application of $\aop$ drives $v$ towards the
  golden ratio, mirroring the triadic process.
  \end{definition}
\end{rsbox}

\begin{rsbox}
  \begin{definition}[Linear Surge]
  The \emph{linear surge} $\lop$ acts on a component $(n,d)$ by
  \begin{equation*}
    \lop(n,d) = (n + d, d).
  \end{equation*}
  This operation adds the denominator to the numerator while leaving the
  denominator unchanged.  It implements the map $v\mapsto v+1$ and yields a
  linear ramp along the mediant lattice.
  \end{definition}
\end{rsbox}

\subsection{Cross–Synchronisation Operator}

\begin{rsbox}
  \begin{definition}[Transformative Reciprocal]
  The \emph{transformative reciprocal} $\psiop$ acts on a state
  $S=(A,B)$ with $A=(n_a,d_a)$ and $B=(n_b,d_b)$ by
  \begin{equation*}
    \psiop(A,B) = ((d_b,n_a),(d_a,n_b)).
  \end{equation*}
  This operator swaps the structural histories encoded in the denominators and
  interchanges numerators and denominators across registers.  It performs a
  discrete quarter–turn in the plane spanned by $(v_A,v_B)$.  Note that
  $\psiop$ is an involution: applying it twice returns the original state.
  \end{definition}
\end{rsbox}

\begin{lemma}[Involution]
The transformative reciprocal satisfies $\psiop^2 = \mathrm{id}$ on the
state space.
\end{lemma}

\begin{proof}
Let $S=(A,B)$ with $A=(n_a,d_a)$ and $B=(n_b,d_b)$.  One application of
$\psiop$ yields $S'=(A',B')$ with $A'=(d_b,n_a)$ and $B'=(d_a,n_b)$.  A
second application produces $(n_b,d_b)$ and $(n_a,d_a)$, which is the original
state.  Thus $\psiop^2$ acts as the identity.
\end{proof}

\begin{remark}[Metaphorical Language]
The denominators carry the cumulative history of operations and can be
considered a form of ``inertia.''  When writing about the action of
$\psiop$, we may occasionally refer to denominators as ``energy'' for
intuition.  This terminology is purely metaphorical and has no ontological
status in RigbySpace.
\end{remark}

% -----------------------------------------------------------------------------
%  Section 3: The Rational Euler Cycle
% -----------------------------------------------------------------------------
\section{The Rational Euler Cycle}

We now assemble the surge and cross–synchronisation operators into a
repeating cycle that emulates a quarter–turn rotation.  This four–step
sequence acts on the state pair and redistributes structural history between
registers.

\begin{rsbox}
  \begin{definition}[Four–Step Euler Cycle]
  Given $S_0=(A_0,B_0)$, one Euler cycle consists of the following steps:
  \begin{enumerate}[label=\arabic*.]
    \item \emph{Real Surge:} $A \leftarrow \aop(A)$.
    \item \emph{Quarter Turn:} $(A,B) \leftarrow \psiop(A,B)$.
    \item \emph{Imaginary Surge:} $B \leftarrow \aop(B)$.
    \item \emph{Quarter Turn:} $(A,B) \leftarrow \psiop(A,B)$.
  \end{enumerate}
  We denote the state after $k$ complete cycles by $S_{4k}$.  Between
  quarter turns, denominators grow via the aggregate surge; the twists
  exchange these growing histories between registers.
  \end{definition}
\end{rsbox}

\begin{lemma}[Denominator Growth]
Let $(n,d)\in \Z_{>0}^2$ with $n > d$.  Applying $\aop$ to $(n,d)$ yields
$(n+d,n)$, whose new denominator $n$ is strictly greater than the original
denominator $d$.  Thus, once a register satisfies $n>d$, its denominator
increases under $\aop$.  Moreover, because denominators are exchanged via
$\psiop$, they do not decrease in subsequent cycles.
\end{lemma}

\begin{proof}
If $(n,d)$ satisfies $n > d$, then $\aop(n,d)=(n+d,n)$ has denominator $n$.
By assumption $n > d$, so the new denominator exceeds the old one.  In
particular, the surge replaces the old denominator with the old numerator.
Since denominators propagate through $\psiop$ without alteration, once the
condition $n > d$ holds for a register, successive applications of $\aop$
maintain strict growth of its denominators.
\end{proof}

\begin{observation}[Damped Harmonic Behaviour]
Let $S_0=(A_0,B_0)$ with $A_0=(n_a,d_a)$, $B_0=(n_b,d_b)$ and
$n_a,d_a,n_b,d_b\in \Z_{>0}$.  Suppose at least one denominator is greater
than one.  Numerical experiments and heuristic reasoning suggest that the
sequence of real register values $v_A(k)$ obtained after each Euler cycle
tends to zero, while the imaginary register values $v_B(k)$ remain bounded
and oscillatory.  In particular, after a finite number of cycles the
denominator of the real register grows beyond its initial numerator, and
subsequent cycles cause $v_A$ to decrease monotonically.  Degenerate cases
where $A_0=B_0=(1,1)$ yield trivial constant behaviour.

\begin{proof}[Heuristic Explanation]
Tracking one cycle shows that the real register evolves from $(n_a,d_a)$ to
$(n_a,n_a+n_b)$, so its value becomes $n_a/(n_a+n_b)$ after one cycle.  The
imaginary register transforms in a complementary fashion.  Because
surge operations cause denominators and numerators to grow, there exists
some cycle index $k$ such that the denominator of the real register
exceeds its numerator.  Beyond this point, each cycle reduces the real
ratio $v_A$ and drives it towards zero.  Meanwhile, the imaginary
ratio $v_B$ oscillates between values determined by the alternating surge
and twist.  A rigorous proof of monotonic decay for all initial
conditions requires establishing explicit recurrences and is left for
future work.  The present analysis views the damped behaviour as an
empirical observation consistent with RS primitives.
\end{proof}
\end{observation}

\begin{remark}[Initial Conditions]
The proposition assumes generic positive integer initial states.  If both
registers start at $(1,1)$, surges leave them unchanged and no dynamics occur.
If one register starts at $(1,1)$ and the other has larger values, the
cycle still produces non‑trivial behaviour because denominators and
numerators are exchanged and grow under successive surges and twists.  The
monotonic decay of $v_A$ is guaranteed once denominators grow beyond
initial numerators.
\end{remark}

\subsection{Qualitative Dynamics}

The rational Euler cycle reproduces the qualitative features of a damped
harmony.  As the real component decays, the imaginary component exhibits
oscillations bounded away from zero and one.  The trajectory spirals towards
a limit cycle in the $(v_A,v_B)$ plane.  Unlike continuous differential
equations, the approach is stepwise; nonetheless, the emergent pattern
mirrors that of a harmonic oscillator with friction.

% -----------------------------------------------------------------------------
%  Section 4: Programmable Rational Synthesizer
% -----------------------------------------------------------------------------
\section{Programmable Rational Synthesizer}

Beyond the canonical Euler cycle, one can engineer discrete waveforms by
altering the sequence of surge and twist operations.  We formalise the notion
of a waveform program and present several examples.

\begin{rsbox}
  \begin{definition}[Waveform Program]
  Let $\{\aop_A, \aop_B, \lop_A, \lop_B, \psiop\}$ denote the set of
  elementary operations acting on the pair $(A,B)$; subscripts indicate which
  register a surge targets.  A \emph{waveform program} is a finite ordered
  list of such operations.  Executing the program sequentially constitutes
  one patch.  Repeating the patch yields a discrete signal extracted from
  the real register $v_A$.
  \end{definition}
\end{rsbox}

We highlight several representative programs:

\begin{enumerate}[label=\arabic*.]
  \item \textbf{Damped Sine Program:} The patch $\{\aop_A, \psiop,
    \aop_B, \psiop\}$ reproduces the Euler cycle.  The real register displays a
    damped oscillation similar to a sine wave.
  \item \textbf{Square Wave Program:} For an integer $m>0$, perform $m$
    aggregate surges on $A$, twist, $m$ surges on $B$, then twist again.
    Charging one register while holding the other fixed and then swapping
    produces a discrete square wave with duty cycle determined by $m$.
  \item \textbf{Sawtooth Program:} For $n>0$, perform $n$ linear surges on
    $A$ followed by a twist.  The real value ramps linearly and is then
    inverted, yielding a sawtooth waveform.
  \item \textbf{Triangle/Trapezoid Program:} For $n>0$, alternate $n$ linear
    surges on $A$ and $n$ on $B$ between twists.  This produces a wave that
    rises linearly, drops sharply, holds while the other register ramps, then
    rises again.
\end{enumerate}

\begin{observation}[Waveform Behaviour]
Each program yields a distinctive discrete signal.  The damped sine program
produces a decaying envelope akin to a harmonic oscillator.  The square wave
program alternates between high and low states.  The sawtooth program
generates a linear ramp followed by a sharp reset.  The triangle/trapezoid
program combines linear rises and plateaus.  These patterns demonstrate the
versatility of RS primitives for synthesising signals without continuous
functions.
\end{observation}

\subsection{External Verification}

RigbySpace restricts itself to integer operations; nonetheless, one can
numerically simulate the above programs to visualise the resulting waveforms.

\begin{verifybox}
  Implementing the surge and twist operators on integer pairs in a short
  script produces sequences of rational values.  Plotting the real register
  over successive patches confirms the qualitative descriptions given above.
  For example, the damped sine program yields a curve that decays towards
  zero while oscillating, and the square wave program toggles between two
  values.  Such plots provide external confirmation but lie outside the
  ontological chain.
\end{verifybox}

% -----------------------------------------------------------------------------
%  Section 5: Conclusion
% -----------------------------------------------------------------------------
\section{Conclusion}

This document demonstrates that rotation and oscillation need not rely on
continuum mathematics.  Through a pair of unreduced rational streams and a
combination of aggregate, linear and cross–synchronisation operations, we
construct a discrete analogue of Euler's rotation.  The rational Euler cycle
produces damped harmonic behaviour, and programmable sequences synthesise
square, sawtooth and triangular waveforms.  The analysis underscores the
expressive power of RS primitives: by carefully orchestrating integer
operations, one can approximate complex behaviours ordinarily attributed to
real or complex analysis.  Future work may explore multi–register couplings
and investigate connections to RS thermodynamics and relativity.

\end{document}